% Chapter 2: Data Presentations and Oracles
% Rewritten against the design-lab Lean 4 kernel (FLT_Proofs.*):
% typed data models, full proofs isomorphic to the verified arguments
% (Text<Informant, VCdim<=Ldim, uniform convergence => PAC), and open problems.
%   time_index, counterexample, advice = Profile A (vocabulary, minimal)
%   noisy_input = Profile A-plus (brief + one theorem reference)
%   iid_sample = Profile G (Framework Bridge: i.i.d. is THE statistical assumption)
%   text_presentation + informant_presentation = Profile D (Separation: text < informant)
%   membership_oracle + equivalence_oracle = Profile G (Framework Bridge: queries vs examples)
%   data_stream = Profile G (Framework Bridge: adversarial vs distributional)

\chapter{Data Presentations and Oracles}
\label{ch:data}

Chapter~\ref{ch:objects} established \emph{what} is learned: concept classes, hypotheses, the objects on which complexity measures act. This chapter asks the question that comes before: \emph{how does data reach the learner?}

The answer cuts deep. A learner given i.i.d.\ samples from an unknown distribution operates in a setting where empirical frequencies converge reliably to true probabilities; the VC dimension decides whether learning is possible. Give the same class to a learner facing an adversarially chosen sequence and the VC dimension becomes silent: thresholds on $\R$ have VC dimension one, yet no algorithm bounds its mistakes against an adversary. Supply positive-only data and the class of all finite languages collapses into the unlearnable. Switch to interactive queries and a class that resists passive learning (deterministic finite automata, under cryptographic hardness) becomes learnable in polynomial query time. Same class, same learner, four outcomes: what changed is the data model.

The data model is the axis along which the field divides into four distinct theories, each with its own characterization theorem, its own complexity measure, and its own impossibility result. Same class, same learner, four outcomes: the variable is the data model.

We begin with vocabulary shared across all paradigms (\S\ref{sec:data-vocab}), then develop each data model in turn, organized by the mathematical content each one carries. Each data model is a type in the companion Lean~4 development, and the three results we prove in full (the text/informant separation, the inequality between the VC and Littlestone dimensions, and the passage from uniform convergence to PAC learnability) are mirrored there.

% ================================================================
% SECTION 1: Shared Vocabulary (Profile A concepts, ~12 lines total)
% ================================================================

\section{Shared Vocabulary}
\label{sec:data-vocab}

Before each paradigm goes its own way, it needs a few shared tools: terms for the mechanical skeleton of any learner-environment interaction. The definitions below are brief because their content is thin; they do real work later as types in the formal development and as building blocks for the richer stream definitions that follow.

A \emph{time index} $t \in \N$ is the discrete step indexing rounds of interaction between learner and environment. A \emph{counterexample} to a hypothesis $h$ with respect to target $c$ is a point $x \in X$ where $h(x) \neq c(x)$, witnessing disagreement; as a type it is the subtype $\{x : h(x) \neq c(x)\}$.\formal{FLT_Proofs.Data.Counterexample} \emph{Advice} is any additional information given to a learner beyond its primary data source; its formal role is explored in the advice reductions of \Cref{ch:exact}.

The most general carrier for sequential data is a \emph{data stream}: a function from time steps to labeled examples, $\text{observe} : \N \to X \times Y$.\formal{FLT_Proofs.Data.DataStream} Texts, informants, and noisy inputs are all data streams with extra structure, defined below as extensions of this type.

\begin{defn}[Noisy Input]
\label{def:noisy-input}
In the \emph{classification noise model} at rate $\eta < 1/2$, each label $y_i$ in the training sample is independently flipped with probability $\eta$: the learner receives $(x_i, y_i \oplus z_i)$ where $z_i \sim \mathrm{Bernoulli}(\eta)$. As a type, a noisy stream extends a data stream with a true-label map and a noise rate bounded by $1/2$.\formal{FLT_Proofs.Data.NoisyDataStream}
\end{defn}

Classification noise does not change the landscape of what is learnable.

\begin{thm}[Noise robustness of PAC learning {\cite{AngluinLaird1988}}]
\label{thm:noise-robust}
For every fixed rate $\eta < 1/2$, a concept class is PAC learnable from clean data if and only if it is PAC learnable under classification noise at rate $\eta$. The sample complexity increases by a factor $O\!\left(1/(1-2\eta)^2\right)$, but the family of learnable classes is identical.
\end{thm}

The forward direction is the substance: a clean learner is recovered from a noisy one by estimating each label's majority vote, the variance of which is controlled by $(1-2\eta)^{-2}$. This robustness is special to the i.i.d.\ model; it fails for adversarial noise in the online setting, which requires different techniques (\cref{op:noise}).

\medskip

With these shared pieces in hand, the four data models can be developed in turn. Each receives treatment proportional to its mathematical content.

% ================================================================
% SECTION 2: The i.i.d. Assumption (Profile G: Framework Bridge)
% ================================================================

\section{The Statistical Data Model: I.I.D.\ Samples}
\label{sec:iid}

\begin{defn}[I.I.D.\ Sample]
\label{def:iid-sample}
An \emph{i.i.d.\ sample} of size $m$ from a distribution $D$ over $X \times Y$ is a sequence $S = ((x_1, y_1), \ldots, (x_m, y_m))$ with each $(x_i, y_i)$ drawn independently from $D$. As a type it bundles a probability measure $D$ on $X \times Y$, a size $m$, and a function $\text{Fin } m \to X \times Y$; the learner knows $m$ but learns about $D$ only through $S$.\formal{FLT_Proofs.Data.IIDSample}
\end{defn}

The i.i.d.\ assumption is the most consequential modeling choice in the field. Everything that makes statistical learning theory work, PAC learnability, uniform convergence, the VC characterization, all of it rests on this single hypothesis about how data arrives. The payoff is uniform convergence; the price is three specific fragilities, each one a door into a different paradigm, and that is why the field has several theories.

\subsection{What i.i.d.\ buys: uniform convergence}

The payoff of the i.i.d.\ assumption is \emph{uniform convergence}: empirical frequencies converge to true probabilities simultaneously across the entire hypothesis class. The word ``uniform'' carries weight here. The required sample size must be fixed before the distribution is chosen and before the target is revealed, which means the same $m_0$ works for every distribution and every target at once. This single-$m_0$ uniformity is what lets a learner pick a hypothesis without knowing the distribution it will be tested against.

\begin{defn}[Uniform convergence]
\label{def:uc}
A hypothesis space $\mathcal{H}$ has \emph{uniform convergence} if for every $\varepsilon, \delta > 0$ there is a sample size $m_0$ such that for \emph{every} distribution $D$ and \emph{every} target $c$, an i.i.d.\ sample of size $m \geq m_0$ satisfies, with probability at least $1 - \delta$,
\[
\sup_{h \in \mathcal{H}} \bigl|\emprisk(h) - \risk(h)\bigr| \leq \varepsilon.
\]\formal{FLT_Proofs.Complexity.Generalization.HasUniformConvergence}
\end{defn}

\begin{thm}[Uniform convergence for VC classes, {\cite{VapnikChervonenkis1971,Valiant1984}}]
\label{thm:uc-vc}
If $\VC(\mathcal{H}) = d < \infty$, then $\mathcal{H}$ has uniform convergence, with $m_0 = O\!\left(\frac{d + \log(1/\delta)}{\varepsilon^2}\right)$ sufficing.\formal{FLT_Proofs.Complexity.Symmetrization.vcdim_finite_imp_uc'}
\end{thm}

The proof, deferred to \Cref{ch:pac}, uses symmetrization and the growth-function bound, both of which exploit i.i.d.\ draws essentially. The payoff is provable here in full: uniform convergence makes empirical risk minimization succeed.

\begin{thm}[Uniform convergence yields PAC learnability]
\label{thm:uc-pac}
If $\mathcal{H}$ is nonempty and has uniform convergence, then $\mathcal{H}$ is PAC learnable, and a consistent learner witnesses it.\formal{FLT_Proofs.Complexity.Generalization.uc_imp_pac}
\end{thm}

\begin{proof}
Define the \emph{consistent learner} $L$: on a sample $S$, output any $h \in \mathcal{H}$ that agrees with every labeled point of $S$ if such an $h$ exists, and otherwise output a fixed $h_0 \in \mathcal{H}$ (available since $\mathcal{H} \neq \varnothing$). Given $\varepsilon, \delta > 0$, let $m_0$ be the sample size from uniform convergence at $(\varepsilon, \delta)$, and set the learner's sample budget to $m_0$.

Fix a target $c \in \mathcal{H}$ and a distribution $D$ realizing it (the labels are $c$'s). The target $c$ is itself consistent with every sample, so a consistent hypothesis always exists, and $L(S)$ has empirical error $\emprisk(L(S)) = 0$. On the uniform-convergence event (which has probability at least $1 - \delta$ for $m \geq m_0$), every hypothesis, $L(S)$ among them, satisfies $|\emprisk - \risk| \leq \varepsilon$. Therefore
\[
\risk(L(S)) \leq \emprisk(L(S)) + \varepsilon = 0 + \varepsilon = \varepsilon
\]
on that event, so $\Pr_{S \sim D^{m}}[\risk(L(S)) \leq \varepsilon] \geq 1 - \delta$. This is exactly the PAC guarantee.
\end{proof}

\Cref{thm:uc-vc,thm:uc-pac} compose into the sufficiency half of the fundamental theorem (\cref{thm:vc-char}): finite VC dimension gives uniform convergence, which gives PAC learnability through ERM. This chain underlies the whole statistical theory, and every link uses the i.i.d.\ assumption.

\subsection{What i.i.d.\ costs: three fragilities}
\label{subsec:iid-costs}

Independence and identical distribution buy uniform convergence. The price is three specific breakdowns, each one a door into a different paradigm.

\begin{enumerate}[leftmargin=*,itemsep=4pt]
\item \textbf{No adversarial robustness.} If an adversary chooses the data sequence, independence fails and uniform convergence no longer implies sequential learnability. The obstruction is a theorem, with thresholds as its witness: there are classes with uniform convergence that are not online learnable (\cref{prop:uc-not-online}). The VC dimension cedes jurisdiction to the Littlestone dimension (\S\ref{sec:streams}).

\item \textbf{No distribution-free sequential guarantees.} The i.i.d.\ model is batch: the learner receives all $m$ examples at once. A sequential version must either fix the distribution across time (the online-with-distributional hybrid, which preserves PAC-style results) or drop distributional assumptions entirely (online learning proper).

\item \textbf{No interaction.} The learner is a passive recipient. It cannot choose which points to query, negotiate with a teacher, or request counterexamples. This passivity is the gap that Angluin's query model fills (\S\ref{sec:queries}).
\end{enumerate}

\begin{prop}[Uniform convergence does not imply online learnability]
\label{prop:uc-not-online}
There is a concept class with uniform convergence (hence finite VC dimension) that is not online learnable. Thresholds on $\R$ are such a class.\formal{FLT_Proofs.Complexity.GeneralizationResults.uc_does_not_imply_online}
\end{prop}

\begin{proof}
Thresholds have $\VC = 1$ (\cref{prop:threshold-vc}), so by \cref{thm:uc-vc} they have uniform convergence. But $\Ldim = \infty$ (\cref{prop:threshold-ldim}): the bisection adversary forces an unbounded number of mistakes against any online learner. By the online characterization (\cref{thm:ldim-char}) infinite Littlestone dimension means not online learnable. Hence uniform convergence does not transport to the adversarial setting.
\end{proof}

\begin{remark}[The VC dimension's jurisdiction]
\label{rmk:vc-jurisdiction}
\Cref{prop:uc-not-online} sharpens a point that is easy to state loosely. The VC dimension characterizes PAC learnability \emph{because of}, and only under, the i.i.d.\ assumption. The same class can have $\VC(\mathcal{C}) = 1$ (PAC learnable from i.i.d.\ data) and $\Ldim(\mathcal{C}) = \infty$ (not online learnable). A complexity measure is not a property of the class alone; it is a property of the class \emph{together with} a data model.
\end{remark}

% ================================================================
% SECTION 3: Text vs Informant (Profile D: Separation Result)
% ================================================================

\section{Enumerative Data: Text and Informant}
\label{sec:text-informant}

Remove the distribution. Remove the sample. Instead, let data arrive as an endless enumeration, one element at a time, with an adversary controlling the order and permitted to withhold any element for as long as it chooses. The learner must eventually commit to a correct hypothesis and stop changing it. This is Gold's model of identification in the limit~\cite{Gold1967}, and it strips away the statistical scaffolding entirely.

Within this stripped setting, a difference that looks quantitative turns out to be qualitative. A \emph{text} shows the learner only positive evidence; an \emph{informant} adds negative evidence as well. One might expect that the second model is just a richer version of the first. The separation proved below shows it is a different paradigm: there is a class identifiable from informant that cannot be identified from any text.

\begin{defn}[Text Presentation]
\label{def:text}
A \emph{text} for a language $L \subseteq \N$ is a data stream $t_1, t_2, \ldots$ of elements of $L \cup \{\#\}$ ($\#$ a pause symbol) in which every element of $L$ appears at least once. As a type, a text for $c : X \to \{0,1\}$ extends a data stream with three conditions: every emitted element is positive, every emitted element satisfies $c$, and every positive point eventually appears.\formal{FLT_Proofs.Data.TextPresentation} A class $\mathcal{C}$ is \emph{identifiable in the limit from text} if some learner, on every $C \in \mathcal{C}$ and every text for $C$, eventually stabilizes on a single correct hypothesis.\formal{FLT_Proofs.Criterion.Gold.EXLearnable}
\end{defn}

\begin{defn}[Informant Presentation]
\label{def:informant}
An \emph{informant} for $L$ is a stream of pairs $(x_i, \ell_i)$ with $\ell_i = \ind[x_i \in L]$ in which every natural number eventually appears. As a type it extends a data stream with two conditions: labels are correct, and every point of $X$ eventually appears.\formal{FLT_Proofs.Data.InformantPresentation}
\end{defn}

A text reports what is in the language; an informant reports membership for every point. The informant gives more, and what it gives is not just quantity: the following theorem shows the gap is qualitative.

\begin{thm}[Text $<$ Informant {\cite{Gold1967}}]
\label{thm:text-informant}
There is a concept class identifiable in the limit from informant but not from text. The class
\[
\mathcal{C}_{\mathrm{fin}} = \{L \subseteq \N : L \text{ finite}\} \cup \{\N\}
\]
of all finite languages together with the full language witnesses the separation.
\end{thm}

\begin{proof}
\emph{From informant: identifiable.} Use the learner that conjectures $\N$ until it sees a negative example, and thereafter conjectures the finite set of positive points seen so far. If the target is $\N$, the informant (being correct) never produces a negative label, so the learner conjectures $\N$ at every step and has converged. If the target is a finite $L$, then $\N \setminus L$ is nonempty and the informant is exhaustive, so a negative example appears at some finite time; from then on the learner conjectures ``positives seen so far,'' and since $L$ is finite, all of its elements appear after finitely many steps, after which the conjecture equals $L$ permanently. Either way the learner stabilizes on a correct hypothesis.

\emph{From text: not identifiable.} Suppose, for contradiction, that a learner $M$ identifies $\mathcal{C}_{\mathrm{fin}}$ from every text. We construct a single text on which $M$ changes its conjecture infinitely often, contradicting convergence. Build the text in stages, maintaining a finite prefix $\sigma$.

Begin enumerating $0, 1, 2, \ldots$, a valid text for $\N$. Since $\N \in \mathcal{C}_{\mathrm{fin}}$ and this is a legitimate text for it, $M$ must eventually conjecture a correct hypothesis for $\N$; let stage $1$ end at the first such time, with prefix $\sigma_1 = \langle 0,1,\ldots,k_1\rangle$.

The prefix $\sigma_1$ is also a legitimate beginning of a text for the finite set $F_1 = \{0,\ldots,k_1\}$. Extend the text by repeating elements of $F_1$ (or padding with $\#$): this makes it a valid text for $F_1 \in \mathcal{C}_{\mathrm{fin}}$, on which $M$ must converge to a hypothesis for $F_1$. Since $F_1 \neq \N$, $M$ must abandon its $\N$-conjecture and conjecture $F_1$ at some later time, ending stage $2$ with prefix $\sigma_2$.

Now resume enumerating fresh natural numbers $k_1+1, k_1+2, \ldots$, again building toward a text for $\N$; $M$ must return to an $\N$-conjecture, ending stage $3$. Iterating, each stage forces $M$ to switch between an $\N$-hypothesis and a finite-set hypothesis. The limiting sequence is a genuine text (every element of its underlying language appears), yet $M$ never stabilizes. This contradicts identification, so no learner succeeds on $\mathcal{C}_{\mathrm{fin}}$ from text.
\end{proof}

The witness $\mathcal{C}_{\mathrm{fin}}$ is instructive: an extremely natural class (all finite sets plus the universal set) already separates the two data models. The obstruction is the impossibility of distinguishing ``all of $\N$'' from ``a large finite set'' using positive data alone.

\begin{remark}[The structural role of negative data]
\label{rmk:negative-data}
The separation shows that negative data does a qualitatively different job: it provides \emph{exclusion}. A text reports what is in $L$; an informant also reports what is \emph{not}, and that single extra power, the ability to say a point is absent, is the whole difference between the two models. Exclusion lets the learner confirm finiteness, which no amount of positive data can do, since any finite prefix of a text for $\N$ is also a prefix of a text for a finite superset. The same asymmetry between inclusion and exclusion recurs in query learning (\S\ref{sec:queries}), where equivalence queries deliver exclusion as counterexamples.
\end{remark}

\begin{historical}
Gold's 1967 paper~\cite{Gold1967} introduced both data models and proved the separation for language identification. It predates PAC learning by 17 years and online learning by 21, making it one of the earliest impossibility results in computational learning theory. The framework remains the standard model for inductive inference and underlies the mind-change hierarchy of Freivalds--Smith and Case--Smith (\Cref{ch:ordinals}), with its criterion hierarchy of finite, explanatory, behaviorally correct, vacillatory, and anomalous identification.\formal{FLT_Proofs.Criterion.Gold.BCLearnable}
\end{historical}

% ================================================================
% SECTION 4: Queries (Profile G: Framework Bridge)
% ================================================================

\section{Query Access: Membership and Equivalence Oracles}
\label{sec:queries}

Both models so far, statistical and enumerative, leave the learner as a passive recipient of data that arrives on someone else's schedule. Angluin's query model~\cite{Angluin1988} breaks that constraint: it grants oracle access, letting the learner ask the questions it needs rather than wait for relevant data to arrive.

\begin{defn}[Membership Oracle]
\label{def:mq}
A \emph{membership oracle} $\mathrm{MQ}$ for a target $c$ answers ``what is $c(x)$?'' for any $x$ chosen by the learner. As a type it bundles a truthful response function with the target it answers about.\formal{FLT_Proofs.Data.MembershipOracle}
\end{defn}

\begin{defn}[Equivalence Oracle]
\label{def:eq}
An \emph{equivalence oracle} $\mathrm{EQ}$ for $c$ answers ``does $h = c$?'': it returns success if $h = c$, and otherwise a counterexample $x$ with $h(x) \neq c(x)$. As a type its query map returns an optional counterexample, with proofs that ``none'' certifies equality and ``some $x$'' is a genuine disagreement.\formal{FLT_Proofs.Data.EquivalenceOracle}
\end{defn}

\begin{defn}[Exact Learning]
\label{def:exact-learning}
A class $\mathcal{C}$ is \emph{exactly learnable} from membership and equivalence queries if some algorithm, for every target $c \in \mathcal{C}$, outputs an $h$ with $h = c$ after a number of queries polynomial in the representation size of $c$ and the size of counterexamples received.\formal{FLT_Proofs.Criterion.PAC.ExactLearnable}
\end{defn}

What these oracles do becomes clear only against what passive data cannot.

\subsection{The power of queries: DFAs as witness}

The comparison is sharpest for DFAs. Passive, distributional data cannot efficiently learn them under standard cryptographic assumptions. With both oracles, Angluin's algorithm learns every $n$-state automaton in polynomial time. The class did not change; the data model did.

\begin{thm}[Angluin's $L^\ast$ algorithm {\cite{Angluin1988}}]
\label{thm:lstar}
The class of $n$-state DFAs over an alphabet $\Sigma$ is exactly learnable using $\mathrm{MQ} + \mathrm{EQ}$ in time polynomial in $n$ and $|\Sigma|$, with at most $n - 1$ equivalence queries and $O(n|\Sigma| m)$ membership queries, where $m$ bounds the length of counterexamples.
\end{thm}

The algorithm builds an observation table encoding the Myhill--Nerode classes of the target, filling it via membership queries and repairing it via equivalence queries; it is developed in full in \Cref{ch:exact}. Contrast the passive setting.

\begin{prop}[DFAs are not PAC-learnable under cryptographic assumptions {\cite{KearnsValiant1994}}]
\label{prop:dfa-pac-hard}
Under standard cryptographic assumptions, DFAs are not efficiently PAC-learnable from random examples alone.
\end{prop}

\begin{remark}[Queries as paradigm shift]
\label{rmk:queries-paradigm}
The DFA example witnesses a genuine paradigm separation: the same class is computationally intractable under passive, distributional data and efficiently learnable under interactive, worst-case queries. Membership queries let the learner \emph{explore} the target's structure; equivalence queries deliver \emph{directed} negative information at the learner's current frontier. The separation runs in both directions: some classes are PAC-learnable but not efficiently exactly learnable (equivalence queries can be costly to simulate), and DFAs go the other way.
\end{remark}

\subsection{Neither oracle suffices alone}

The combination $\mathrm{MQ} + \mathrm{EQ}$ is necessary because each oracle has a specific weakness the other supplies. Membership queries explore structure but cannot certify global correctness; equivalence queries certify correctness but cannot probe structure efficiently. The two propositions below make each weakness precise.

\begin{prop}[Membership queries alone cannot verify]
\label{prop:mq-alone}
Let $X$ be infinite and let $\mathcal{C}$ contain two concepts $c_1 \neq c_2$. A learner restricted to membership queries cannot exactly identify the target after any finite number of queries, unless its queries already pin down the target on all of $X$.
\end{prop}

\begin{proof}
After $q$ membership queries the learner has observed the target on a finite set $Q$ of points. If $c_1$ and $c_2$ agree on $Q$ but differ at some $x \notin Q$ (which can be arranged whenever the queried points fail to separate the class), then the transcript is identical whether the target is $c_1$ or $c_2$, so any fixed output is wrong for one of them. Membership queries supply local values but no certificate of global correctness; only an equivalence query can certify $h = c$.
\end{proof}

\begin{prop}[Equivalence queries alone learn finite classes, but slowly]
\label{prop:eq-alone}
A finite class $\mathcal{C} = \{c_1, \ldots, c_N\}$ is exactly learnable from equivalence queries alone using at most $N - 1$ queries; in the worst case $N - 1$ are necessary, which is exponential in the representation size.
\end{prop}
\begin{proof}
Maintain the version space of candidates consistent with all counterexamples received, initialized to $\mathcal{C}$. Query any surviving candidate. A success ends the process; a counterexample $x$ eliminates at least the queried candidate (it disagrees with the target at $x$) and is consistent only with the candidates that survive. Each query removes at least one candidate, so after at most $N - 1$ queries a single candidate remains and is correct. An adversary answering against the largest surviving set forces one elimination per query, giving the matching lower bound. Equivalence queries certify globally but carry no mechanism to probe structure, so the count is governed by $|\mathcal{C}|$ rather than by representation size.
\end{proof}

% ================================================================
% SECTION 5: Adversarial Streams (Profile G: Framework Bridge)
% ================================================================

\section{Adversarial Streams and the Online Data Model}
\label{sec:streams}

\begin{defn}[Data Stream (Online Protocol)]
\label{def:data-stream}
In the \emph{online protocol}, data arrives as an adversarially chosen stream. At each round $t = 1, 2, \ldots$: the adversary selects $x_t \in X$ (possibly depending on the learner's past predictions); the learner predicts $\hat{y}_t$; the true label $y_t = c(x_t)$ is revealed. The learner minimizes the total mistake count $|\{t : \hat{y}_t \neq y_t\}|$, and there is \emph{no distributional assumption}. A class is \emph{online learnable} if some learner has a finite worst-case mistake bound over all targets in the class and all adversarial sequences.\formal{FLT_Proofs.Learner.Core.OnlineLearner}
\end{defn}

Strip away everything the i.i.d.\ model provides: the distribution, the independence, the convergence of frequencies. What remains is a sequential game: the adversary places instances one by one, adapting to the learner's past predictions; the learner must predict each label before seeing it; the running score is mistakes. No distributional guarantee softens the adversary's choices. This austerity is not just philosophical severity; it forces a different complexity measure, one the VC dimension cannot see.

\begin{defn}[Mistake tree and Littlestone dimension]
\label{def:mistake-tree-data}
A \emph{mistake tree} of depth $d$ is a complete binary tree whose internal nodes are labeled by instances $x \in X$; the inductive type is a leaf or a branch carrying an instance and two subtrees.\formal{FLT_Proofs.Complexity.GameInfra.LTree} The tree is \emph{shattered} by $\mathcal{C}$ if, recursively, the root instance $x$ admits a concept in $\mathcal{C}$ labeling it $1$ whose restriction shatters the left subtree and a concept labeling it $0$ whose restriction shatters the right subtree, so that every root-to-leaf path is realized by some concept forcing a mistake at each node.\formal{FLT_Proofs.Complexity.GameInfra.LTree.isShattered} The \emph{Littlestone dimension} $\Ldim(\mathcal{C})$ is the greatest depth of a shattered mistake tree.\formal{FLT_Proofs.Complexity.GameInfra.LittlestoneDim}
\end{defn}

\begin{warning}[Two shattering conventions]
\label{rmk:branchwise}
There are two readings of ``shattered.'' The \emph{branch-wise} reading asks only that each node admit some concept of the right label, with no consistency required along a path; the \emph{path-consistent} reading (above, and used in the characterization theorem) restricts the class at each branch and demands one concept realize an entire path. The branch-wise reading over-counts: the class $\{x \mapsto 1, x \mapsto 0\}$ of two constant concepts has branch-wise dimension $\infty$ yet is online learnable with a single mistake. The path-consistent dimension is the correct invariant; we use it throughout.
\end{warning}

The Littlestone dimension is to online learning what the VC dimension is to PAC learning, and the two are ordered.

\begin{prop}[The VC dimension lower-bounds the Littlestone dimension]
\label{prop:vc-le-ldim}
For every concept class, $\VC(\mathcal{C}) \leq \Ldim(\mathcal{C})$.\formal{FLT_Proofs.Complexity.Littlestone.BranchWiseLittlestoneDim_ge_VCDim}
\end{prop}

\begin{proof}
Let $S = \{x_1, \ldots, x_d\}$ be shattered by $\mathcal{C}$ in the VC sense, with $d = \VC(\mathcal{C})$. Build the complete binary mistake tree $T$ of depth $d$ that queries $x_1$ at the root and, along every branch, queries $x_2, \ldots, x_d$ in order (the same instance at every node of a given level). Each root-to-leaf path corresponds to a labeling $b \in \{0,1\}^d$ of $S$. Because $S$ is shattered, some concept $c \in \mathcal{C}$ realizes exactly the labeling $b$ on $S$, and this single concept forces the prescribed mistake at every node of the path; this is precisely path-consistent shattering. Hence $T$ is a shattered mistake tree of depth $d$, so $\Ldim(\mathcal{C}) \geq d = \VC(\mathcal{C})$.
\end{proof}

The inequality can be arbitrarily loose. Thresholds on $\R$ have $\VC = 1$ but $\Ldim = \infty$ (\cref{prop:threshold-vc,prop:threshold-ldim}), so finite VC dimension does not bound the Littlestone dimension, which is the formal content of fragility~(1) above. What the Littlestone dimension buys in return is a clean characterization and an algorithm that meets it.

\begin{thm}[Littlestone characterization {\cite{Littlestone1988}}]
\label{thm:ldim-char}
A concept class is online learnable if and only if $\Ldim(\mathcal{C}) < \infty$, and the optimal worst-case mistake bound equals $\Ldim(\mathcal{C})$, achieved by the Standard Optimal Algorithm.\formal{FLT_Proofs.Theorem.Online.littlestone_characterization}
\end{thm}

The Standard Optimal Algorithm predicts, at each point, the label whose surviving version space has the larger Littlestone dimension; each mistake strictly decreases that dimension, so at most $\Ldim(\mathcal{C})$ mistakes occur. The characterization and the matching bound are developed in \Cref{ch:online}.\formal{FLT_Proofs.Theorem.Online.optimal_mistake_bound_eq_ldim} Online learnability is strictly stronger than PAC learnability: every online-learnable class is PAC learnable, and thresholds show the converse fails.\formal{FLT_Proofs.Theorem.Separation.online_strictly_stronger_pac}

\begin{remark}[What adversarial data buys: robustness]
\label{rmk:online-robustness}
A mistake-bounded online learner guarantees its bound against \emph{any} sequence: worst-case, non-stationary, distribution-shifting. The bound $\Ldim(\mathcal{C})$ holds regardless of how data is generated. This unconditional guarantee is what makes online learning the right framework when distributional assumptions are unjustified, and it is exactly the strength that the i.i.d.\ model trades away for the smaller invariant $\VC \leq \Ldim$.
\end{remark}

% ================================================================
% SECTION 6: Taxonomy (all four models compared)
% ================================================================

\section{Taxonomy: Four Data Models, Four Theories}
\label{sec:taxonomy}

The four data models have now been developed individually. Placed side by side, they form a definite structure: each is characterized by a distinct combinatorial dimension, the paradigms are ordered by strict containment in some directions and separated by verified impossibility in others, and every claim in Figure~\ref{fig:data-taxonomy} is backed by a machine-checked proof.\formal{FLT_Proofs.Theorem.Separation.online_strictly_stronger_pac}

\begin{figure}[ht]
\centering
\begin{tikzpicture}[
    cell/.style={draw, thick, rounded corners, minimum width=2.9cm, minimum height=1.05cm, font=\scriptsize, align=center},
    contain/.style={-{Stealth[length=2mm]}, semithick, defblue},
    strict/.style={-{Stealth[length=2mm]}, semithick, thmred},
    lit/.style={{Stealth[length=1.5mm]}-{Stealth[length=1.5mm]}, dashed, notegray, semithick},
    litd/.style={-{Stealth[length=1.5mm]}, dashed, notegray, semithick},
    elab/.style={font=\tiny, text=notegray, inner sep=1.3pt, fill=white, align=center},
    slab/.style={font=\tiny, text=thmred!80!black, inner sep=1.3pt, fill=white, align=center}
]
% compound nodes: data model + its characterizing dimension
\node[cell, fill=onlineblue!12]         (str)  at (0, 2.85)   {\textbf{adversarial stream}\\Littlestone dim};
\node[cell, fill=defblue!12]            (iid)  at (0, 0.7)    {\textbf{i.i.d.\ sample}\\VC dim};
\node[cell, fill=sepgreen!12]           (text) at (-3.55,-1.5){\textbf{enumerative text}\\mind-change ordinal};
\node[cell, fill=edgeorange!12, dashed] (mq)   at (3.55,-1.5) {\textbf{MQ\,+\,EQ}\\query complexity};
\node[cell, fill=layergray, dashed, minimum height=0.8cm] (ctx) at (0,-3.35) {\textbf{in-context (attention)}};
% spine: verified containment + strict separation (online vs PAC)
\draw[contain] ([xshift=-7pt]str.south) -- node[elab,left]{online $\subseteq$ PAC} ([xshift=-7pt]iid.north);
\draw[strict]  ([xshift=7pt]iid.north) -- node[slab,right]{$\not\Rightarrow$ thresholds:\\VC$=1$, Ldim$=\infty$} ([xshift=7pt]str.south);
% Gold vs PAC (verified strict)
\draw[strict] (text) -- node[slab,below right,pos=0.6]{$\not\Rightarrow$ finite support} (iid);
% literature / open edges
\draw[litd] (mq) -- node[elab,above right,pos=0.5]{exact $\not\Rightarrow$ PAC\\(DFAs; crypto, lit.)\\reverse open: \cref{op:reverse-separation}} (iid);
\draw[litd] (ctx) -- node[elab,right]{measurability\\side-condition} (iid);
\node[font=\tiny, text=sepgreen!55!black, anchor=north, inner sep=1pt] at (text.south) {\leanname{gold_theorem}: text obstruction};
% legend
\node[font=\tiny, anchor=west, text=notegray] at (-5.1,-4.35) {%
  \textcolor{defblue}{\rule[0.45ex]{0.5cm}{0.9pt}}~containment (verified) \quad
  \textcolor{thmred}{\rule[0.45ex]{0.5cm}{0.9pt}}~$\not\Rightarrow$ strict separation (verified) \quad
  \textcolor{notegray}{\rule[0.45ex]{0.15cm}{0.8pt}\,\rule[0.45ex]{0.15cm}{0.8pt}}~literature / open};
\end{tikzpicture}
\caption[The four data models, their characterizing dimensions, and the proven separations]{Each node pairs a data model with the combinatorial quantity that decides learnability within it: the VC dimension for i.i.d.\ samples, the Littlestone dimension for adversarial streams, the mind-change ordinal for enumerative texts, query complexity for oracle access. The edges prove what the nodes cannot say alone. The red arrows are verified separations: the same class, thresholds on $\R$, has VC dimension one yet Littlestone dimension infinite, so finite VC dimension does not carry over to the adversarial setting. This is the chapter's argument in diagram form: a complexity measure is not a property of the class in isolation; it is a property of the class under a specific data model.}
\label{fig:data-taxonomy}
\end{figure}

\Cref{tab:data-comparison} displays the correspondence between data model, complexity measure, characterization theorem, and impossibility result.

\begin{table}[ht]
\centering
\small
\caption{Each row is a complete paradigm: the data model fixes the complexity measure, which fixes the characterization theorem, which fixes the impossibility result. Change the data model and every column changes with it.}
\label{tab:data-comparison}
\begin{tabularx}{\textwidth}{>{\bfseries\small}l X X X X}
\toprule
& \textbf{Data Model} & \textbf{Complexity Measure} & \textbf{Characterization} & \textbf{Key Impossibility} \\
\midrule
PAC &
I.I.D.\ sample from unknown $D$ &
$\VC(\mathcal{C})$ &
$\mathcal{C}$ is PAC-learnable iff $\VC(\mathcal{C}) < \infty$ \cite{Valiant1984} &
NFL: $\{0,1\}^X$ not learnable \\[6pt]

Gold &
Text (positive) or informant (positive + negative) &
Finite thickness, mind-change ordinals &
Angluin's condition; telltale sets \cite{Gold1967} &
$\mathcal{C}_{\mathrm{fin}} \cup \{\N\}$: text $<$ informant \\[6pt]

Exact &
MQ + EQ (interactive) &
Query complexity (polynomial in representation) &
$L^\ast$ for DFAs; class-specific \cite{Angluin1988} &
MQ alone or EQ alone insufficient \\[6pt]

Online &
Adversarial stream (no distribution) &
$\Ldim(\mathcal{C})$ &
$\mathcal{C}$ is online-learnable iff $\Ldim(\mathcal{C}) < \infty$ \cite{Littlestone1988} &
$\VC < \infty \not\Rightarrow \Ldim < \infty$ \\
\bottomrule
\end{tabularx}
\end{table}

\section{Cross-Paradigm Separations}
\label{sec:data-separations}

The three separations below are the chapter's principal content. Each one names a class that is learnable under one data model and not another, together with the explicit mechanism blocking transfer. Each has a verified proof in the Lean~4 development, developed in full in \Cref{ch:separations}.

\begin{center}\small
\begin{tabular}{llll}
\toprule
\textbf{Separation} & \textbf{Witness} & \textbf{Mechanism} & \textbf{Reference} \\
\midrule
Text $<$ Informant & $\mathcal{C}_{\mathrm{fin}} \cup \{\N\}$ & Diagonalization & \cite{Gold1967} \\[3pt]
PAC $\not\Rightarrow$ Online & Thresholds on $\R$ & $\VC = 1$, $\Ldim = \infty$ & \cite{Littlestone1988} \\[3pt]
Passive $\not\Rightarrow$ Exact & DFAs & Crypto.\ hardness of PAC; & \cite{Angluin1988}, \\
(and vice versa) & & $L^\ast$ in poly queries & \cite{KearnsValiant1994} \\
\bottomrule
\end{tabular}
\end{center}

Each separation has the same logical structure: a witness class learnable under one data model but not another, together with an explicit obstruction. The PAC/online separation is verified directly (some class is PAC learnable but not online learnable\formal{FLT_Proofs.Theorem.Separation.pac_not_implies_online}) and extends to a three-way separation of PAC, online, and Gold identification.\formal{FLT_Proofs.Theorem.Separation.online_pac_gold_separation} The witnesses (thresholds, finite sets, finite automata) are among the simplest classes in the theory. Their simplicity is the signature of a genuine paradigm boundary: even the simplest classes distinguish the paradigms.

\section{What This Chapter Established}
\label{sec:ch2-summary}

The chapter's argument is one sentence: the data model determines the theory. More precisely:

\begin{enumerate}[leftmargin=*,itemsep=3pt]
\item \textbf{Four data models yield four paradigms.} I.i.d.\ samples give PAC learning; enumerative texts give Gold identification; query oracles give exact learning; adversarial streams give online learning. These are different frameworks with different characterization theorems, complexity measures, and impossibility results, and each data model is a distinct type.

\item \textbf{The separations are witnessed by simple classes.} Thresholds separate PAC from online (\cref{prop:vc-le-ldim} and the threshold gap); $\mathcal{C}_{\mathrm{fin}} \cup \{\N\}$ separates text from informant (\cref{thm:text-informant}); DFAs separate passive from interactive. The simplicity of the witnesses shows the boundaries are fundamental.

\item \textbf{The complexity measure is a joint property of class and data model.} The VC dimension measures $\mathcal{C}$ against i.i.d.\ data; the Littlestone dimension measures it against adversarial data; $\VC \leq \Ldim$ orders them, and the gap is unbounded. Neither is ``the'' complexity of $\mathcal{C}$; each is the complexity of the class under a presentation. The same object has different complexities under different data models; this is a central structural insight of the field, and the reason uniform convergence does not transport to the adversarial setting (\cref{prop:uc-not-online}).
\end{enumerate}

Chapters~\ref{ch:pac}--\ref{ch:universal} develop each paradigm in full. The present chapter has identified the axis along which they separate.

% ================================================================
% TRANSFORMER CONNECTION (retrofit): attention as a measurable batch
% learner over i.i.d. data; routing as a measurable partition;
% the float32 forward-pass error of the readout.
% ================================================================

\section{The Data of a Transformer}
\label{sec:ch2-transformer}

The consistent learner of \cref{thm:uc-pac} was defined inline: ``output any $h \in \mathcal{H}$ agreeing with $S$,'' with no $\sigma$-algebra on $\mathcal{H}$ and no requirement that $S \mapsto L(S)$ be measurable. For the simplest classes that omission is harmless. For a parametric attention model it is not free, and making it explicit reveals two axes the classical model suppresses: the \emph{measurability} of the data-to-hypothesis map, and the \emph{finite precision} in which data is actually read by a real machine.

\begin{prop}[Attention routing is a measurable partition]
\label{prop:attention-measurable-routing}
The data-dependent routing of a multi-head attention model (the partition of the query space according to which key each query attends to) is measurable, and the hypothesis class it induces is a well-behaved VC-measurability target.\formal{TLT_Proofs.Routing.MeasurableScoreRouting.multiHeadArgmax_wellBehaved}  The same guarantee holds for a finite-cell router.\formal{TLT_Proofs.Tame.FiniteCellRouter.finiteCellRouter_wellBehaved}  Consequently the attention learner is a \emph{measurable} batch learner in the sense of the condition \leanname{LearnEvalMeasurable}\formal{FLT_Proofs.Theorem.Separation.LearnEvalMeasurable} that \cref{thm:uc-pac} left implicit.
\end{prop}

\begin{proof}
Fix the parameters and a query~$x$.  The routing decision is the $\arg\max$ over keys of the scores $\langle x, k_i\rangle$, so the set of queries routed to key~$i$ is the intersection of the half-spaces $\{x : \langle x, k_i - k_j\rangle \geq 0\}$ over $j \neq i$, which is a Borel set.  The routing map $(\text{parameters}, x) \mapsto (\text{winning key})$ is therefore measurable, and the head's prediction is a measurable function of parameters and query.  A parametrized family of concepts with this joint measurability is by definition a well-behaved VC-measurability target, and the learner's selection $S \mapsto L(S)$ inherits measurability from it.  This is the condition \leanname{LearnEvalMeasurable}: for every sample size~$m$, the evaluation $(S, x) \mapsto L(S)(x)$ is measurable, which is what the uniform-convergence argument of \cref{thm:uc-pac} requires.
\end{proof}

The second axis is invisible in the real-valued data model of \cref{def:iid-sample} and unavoidable in practice: the learner does not read real numbers.

\begin{prop}[The forward pass reads data in finite precision]
\label{prop:attention-fp32}
On real hardware the attention readout is a \texttt{float32} fold-sum of softmax-weighted values, the idealized real-valued sum being unavailable to a finite-precision machine.  The executed result differs from the ideal by a bounded amount: for a fold of $m$ summands whose absolute values total at most~$S$,
\[
  \bigl|\,\mathrm{fp32fold}(v) - \textstyle\sum_i v_i\,\bigr| \;\leq\; 2^{-24}\,(m+1)\,S,
\]\formal{TLT_Proofs.Bridge.Fp32.ClampedReductionRounding.fp32Foldl_error_le_of_sum_bound}
and the executed IEEE-754 fold obeys an analogous accumulated-error envelope.\formal{TLT_Proofs.Bridge.Fp32.SequentialSummationBackwardError.ie32_foldl_envelope}
\end{prop}

\begin{proof}[Proof idea]
Each addition in the fold commits a relative rounding error of at most the unit roundoff $2^{-24}$; propagating these through the $m$ summations and bounding by the total magnitude~$S$ gives the stated linear-in-$m$ envelope by backward-error analysis.  What matters is that such a constant exists at all: the map from data to prediction that actually runs is the idealized one perturbed by a bounded amount.
\end{proof}

\begin{remark}[A third coordinate of the data model]
\label{rmk:attention-data-model}
\Cref{rmk:vc-jurisdiction} held that a complexity measure is a property of the class \emph{together with} a data model.  Attention adds a coordinate that the four classical models share but never expose: the \emph{precision} in which the data is represented.  The i.i.d.\ jurisdiction is where the measurable routing of \cref{prop:attention-measurable-routing} lets the VC machinery apply at all; the float32 envelope of \cref{prop:attention-fp32} is the cost of executing that machinery on data a real machine can hold.  The classical theory idealizes both axes away: it assumes the learner measurable and the arithmetic exact, and a transformer is where neither assumption is free.
\end{remark}

\section{Exercises and Open Problems}
\label{sec:ch2-open}

These problems sit at the edges where the chapter's argument is not yet complete. Each is anchored to a specific construct above. The tags signal what kind of incompleteness you face: a \emph{known unknown} (\textsc{ku}) is a question whose answer is open; an \emph{unknown known} (\textsc{uk}) is a result the literature already contains, but one not yet formalized or absorbed into a clean general statement in this framework. In both cases the question is real, the difficulty is genuine, and the answer, if you find one, changes something.

\begin{openprob}[Formalizing noise robustness, \textsc{uk}]
\label{op:noise}
Formalize \cref{thm:noise-robust} (Angluin--Laird), the equivalence ``PAC under classification noise at rate $\eta < 1/2$ $\Leftrightarrow$ clean PAC'' with the explicit $O(1/(1-2\eta)^2)$ overhead, and identify the exact threshold behavior as $\eta \to 1/2$.
\end{openprob}

\begin{openprob}[The optimal uniform-convergence constant, \textsc{ku}]
\label{op:uc-constant}
\Cref{thm:uc-vc} gives $m_0 = O((d + \log(1/\delta))/\varepsilon^2)$. Determine the exact leading constant for VC classes: symmetrization, chaining, and the Dudley entropy integral give different bounds, and which is tight for general classes of dimension $d$ is open.
\end{openprob}

\begin{openprob}[A characterization of identifiability from text, \textsc{uk}]
\label{op:text-char}
\Cref{thm:text-informant} shows $\mathcal{C}_{\mathrm{fin}} \cup \{\N\}$ is not text-identifiable. Angluin's telltale condition characterizes which classes are: a class is text-identifiable iff every language has a finite ``telltale'' subset distinguishing it from its sublanguages in the class. Formalize this characterization on top of the typed criterion hierarchy (\leanname{FLT_Proofs.Criterion.Gold.EXLearnable}).
\end{openprob}

\begin{openprob}[Closing the VC--Littlestone gap, \textsc{ku}]
\label{op:vc-ldim-gap}
\Cref{prop:vc-le-ldim} gives $\VC \leq \Ldim$, and thresholds make the gap infinite. Characterize the classes for which $\Ldim$ is finite given finite $\VC$ (the Littlestone classes), and give a quantitative bound on $\Ldim$ in terms of $\VC$ and a second structural parameter for the cases where one exists. When is $\Ldim = \VC$ exactly?
\end{openprob}

\begin{openprob}[Query complexity of exact learning, \textsc{ku}]
\label{op:query-complexity}
For the DFA learner of \cref{thm:lstar}, the equivalence-query count $n-1$ is optimal but the membership-query count $O(n|\Sigma| m)$ is not known to be. Determine the optimal joint $\mathrm{MQ} + \mathrm{EQ}$ complexity of exactly learning $n$-state DFAs, and characterize the classes for which $\mathrm{EQ}$ alone (no membership queries) suffices in a polynomial number of queries.
\end{openprob}

\begin{openprob}[The minimally adequate teacher, \textsc{uk}]
\label{op:mat}
\Cref{prop:mq-alone,prop:eq-alone} show each oracle alone is insufficient. Characterize exactly which concept classes become efficiently exactly learnable when the two oracles are combined (the ``minimally adequate teacher'' question) and relate this to the VC and Littlestone dimensions, since neither passive invariant governs query learning.
\end{openprob}

\begin{openprob}[An invariant for query learning, \textsc{ku}]
\label{op:query-dimension}
PAC learning has the VC dimension and online learning has the Littlestone dimension, each characterizing its paradigm. Exact learning has no single characterizing invariant. Construct a combinatorial dimension that characterizes polynomial-query exact learnability the way $\VC$ and $\Ldim$ characterize their paradigms, or prove that no such single invariant exists.
\end{openprob}

\begin{openprob}[Distribution shift and online-to-batch, \textsc{ku}]
\label{op:online-batch}
\Cref{rmk:online-robustness} notes that online bounds hold under distribution shift while i.i.d.\ guarantees do not. Quantify the degradation: given a class of Littlestone dimension $\ell$ and a sequence whose distribution drifts with total variation budget $B$, give matched upper and lower bounds on the achievable risk, interpolating between the i.i.d.\ ($B = 0$) and adversarial regimes.
\end{openprob}

\begin{openprob}[A unified data-model complexity, \textsc{uk}]
\label{op:unified-measure}
The four paradigms use four invariants (VC, mind-change ordinal, query complexity, Littlestone). Determine whether they are instances of a single complexity functor on (class, data-model) pairs, recovering each invariant by specializing the data model, and classify which pairs of data models admit a transfer theorem versus a provable separation (\cref{tab:data-comparison}).
\end{openprob}

\begin{openprob}[Verifying the three-way separation constructively, \textsc{uk}]
\label{op:three-way}
The PAC/online/Gold separation is verified (\leanname{FLT_Proofs.Theorem.Separation.online_pac_gold_separation}), but the witnesses are produced existentially. Give a constructive, effective version: a single family of classes, indexed by a parameter, that realizes each pairwise separation with explicit witnesses and explicit obstructions, so that the entire separation lattice of \cref{tab:data-comparison} is exhibited by one parameterized construction rather than three independent ones.
\end{openprob}

\begin{openprob}[The measurable attention learner, \textsc{uk}]
\label{op:attention-measurable}
\Cref{prop:attention-measurable-routing} shows the attention learner is measurable, the condition \leanname{LearnEvalMeasurable} that \cref{thm:uc-pac} left implicit, with the routing well-behavedness as witness (\leanname{multiHeadArgmax_wellBehaved}).  Assemble these into a single statement: a parametric softmax-attention learner over i.i.d.\ data is PAC, with the measurability discharged rather than assumed.
\end{openprob}

\begin{openprob}[Does PAC survive finite precision? \textsc{ku}]
\label{op:attention-fp32-pac}
\Cref{prop:attention-fp32} bounds the gap between the executed float32 attention learner and its idealized real-valued counterpart (\leanname{fp32Foldl_error_le_of_sum_bound}).  Does the PAC guarantee transfer to the executed learner, and at what cost: is there a sample-complexity penalty that scales with the precision (the unit roundoff $2^{-24}$), so that learnability degrades gracefully as arithmetic coarsens?
\end{openprob}

\begin{openprob}[Attention under the other data models, \textsc{ku}]
\label{op:attention-data-models}
The mapping of \cref{sec:ch2-transformer} is i.i.d.\ (batch).  Where does attention sit under the other three data models of this chapter: is there an online (streaming) attention learner with a finite mistake bound, or a query-driven (active) one, and does the measurable-routing well-behavedness (\leanname{finiteCellRouter_wellBehaved}) transfer to those presentations, or does the data model change attention's learnability as it changes thresholds'?
\end{openprob}

\subsection*{Counterfactual exercises}

The chapter's argument is that the data model determines everything downstream: the complexity measure, the characterization, the impossibility. The exercises below press on that claim directly. Each one takes a single feature of a data model, corrupts it, weakens it, or swaps it for an adversarial counterpart, and asks what breaks. The answer in each case either confirms the paradigm boundary or locates the exact condition that holds the theory together.

\begin{enumerate}[leftmargin=*,itemsep=3pt]
\item \textbf{Noisy oracle.} Corrupt labels at a rate below one half. Show that the VC dimension is invariant under the noise, so PAC learnability is preserved (\cref{thm:noise-robust}).
\item \textbf{Semi-supervised text.} Add a few labels to a positive-only text. Pin down the condition on the added labels under which the class becomes identifiable, that is, when they supply the exclusion information the finite-support witness needs.
\item \textbf{Stream with bounded memory.} Cap the learner's state. Explain why this can only raise the mistake bound above the Littlestone dimension, and why no characterizing dimension is known for the memory-bounded model.
\item \textbf{Distribution shift.} Let the data be non-stationary. Show that the Littlestone mistake bound covers a drifting distribution verbatim, since it holds against any sequence (\leanname{optimal_mistake_bound_eq_ldim}).
\item \textbf{Partial labels.} Let the oracle abstain. Identify the abstention rate at which the surviving labels recover enough exclusion to defeat the Gold diagonalization (\leanname{gold_theorem}), tipping the model from text-like to informant-like.
\item \textbf{Adversarial versus stochastic order.} Switch the sample order from random to adversarial. Stochastic order makes VC govern, adversarial order makes Littlestone govern; name the witness (thresholds, VC one, Littlestone infinite) that forces the gap (\leanname{online_strictly_stronger_pac}).
\item \textbf{Privacy constraint.} Require a differentially private learner. Argue that under pure privacy the governing quantity is a representation dimension distinct from VC, so the i.i.d.\ characterization no longer applies.
\item \textbf{In-context presentation.} Let the data arrive as a prompt consumed by an attention learner. A hard attention head's failure event can be analytic while failing to be Borel (\leanname{attention_nonBorel_badEvent}); exhibit the measurability side-condition under which the VC characterization still applies.
\end{enumerate}
